\chapter{A World Is a Field}

\section{The Difference Between a Picture and a Place}

Chapter 1 began with a simple observation: a player returns to the same stone wall in a virtual world and finds its cracks and moss subtly changed. At that point the example served only as a symptom, one instance among several of a failure this book had not yet named precisely. This book now has enough machinery in place, or is about to, to ask why the symptom occurs and what kind of architecture would make it impossible. The wall is the same wall. What has changed is the question being asked of it.

Return, then, to the experiment in its more deliberate form. Look at a stone wall, turn away, and look back. In the physical world this act is uneventful: the wall has not changed, and any difference in what you see is attributable to lighting, your position, or the passage of time acting on the stone itself. The wall persisted while you were not looking at it, and your second glance is a second observation of the same underlying thing.

Now perform the analogous experiment inside a generative system. Ask a model to render a stone wall. Turn away, in the sense of ending the generation. Ask again. The model produces a second wall: plausible, perhaps beautiful, quite possibly indistinguishable in quality from the first. But nothing in the system guarantees that the two renderings depict the same wall. The stones may be arranged differently. The moss may have migrated. The crack that ran diagonally through the third course of stones may simply not be there the second time, not because it healed, but because it was never a fact about anything. It was a detail invented once, for one rendering, and then forgotten completely.

This is easy to dismiss as a problem of visual consistency, the kind of thing a texture cache or a fixed random seed might patch over. That diagnosis is too shallow. The issue is not that the two walls look different. The issue is that there is no wall. There is only the act of generating a wall, performed twice, with no structure connecting the two acts. Consistency, when it happens, is a statistical accident of the model's priors rather than a consequence of anything persisting between the two renderings. A system built this way has not built a place. It has built a very capable picture generator, and a place is not a sequence of pictures, however good the pictures are.

The distinction matters because it is ontological rather than aesthetic. An observer's confidence that they are looking at the same wall twice depends on a claim about what exists between observations, not merely a claim about what any single observation looks like. Current generative systems, from diffusion models to autoregressive world models operating on tokens or frames, are built almost entirely around the second kind of claim. They are extraordinarily good at producing an observation that is locally plausible. They contain no mechanism, and typically no representation at all, for the thing an observation is supposed to be an observation of. This chapter is about supplying that missing thing.

\section{Reality Requires Persistence}

The claim at the center of this book can be stated plainly before any mathematics is introduced:

> A world is not the sequence of its observations. A world is the persistent structure from which observations arise.

This sentence reverses the implicit assumption behind most generative architectures. Those architectures treat the observation as primary: the frame, the token, the rendered image is the object of interest, and whatever comes before it exists only to make the next observation plausible given the previous ones. Context windows, latent caches, and conditioning histories are all, in the end, mechanisms for making observation number *n* consistent with observations *1* through *n − 1*. They are bookkeeping devices layered on top of a fundamentally observation-first ontology.

The alternative treats the observation as derived. There is a persistent entity, call it the world-state, and what an agent sees at any moment is a partial, situated projection of that state rather than the state itself. Formally, we will eventually write this as an observation operator,

\[ O_\theta(M_t) = y_t, \]

where \(M_t\) denotes the world-state at time t and \(y_t\) denotes what an observer actually perceives. The subscript θ marks that this projection is itself a model, parameterized and learned, standing in for a camera, a sensor, a rendering pipeline, or a language description depending on the modality in question. The crucial point, for now, is not the precise form of \(O_\theta\) but its role. It is a *reading* operation performed on something that already exists. It does not manufacture the wall. It reports on the wall.

This single reversal has consequences that ripple through everything that follows. If observation is derived from state, then two observations of the same region should agree not because a model was trained to make them agree, but because they are readings of the same underlying structure, taken at nearby moments, through a stable operator. Disagreement becomes diagnostic rather than expected. A crack that appears and disappears is no longer an unremarkable feature of sampling variance; it is evidence that the system has no persistent state to be unfaithful to, or that the operator reading that state is unstable, or that the state itself was corrupted between readings. Once persistence is taken seriously as a requirement, inconsistency stops being invisible noise and becomes a signal that something specific has gone wrong, and can in principle be traced to a specific failure in a specific part of the architecture.

Note what this section has and has not done. It has argued that observation must be derived from a persistent state, and it has named the operator that performs that derivation. It has not yet said what that persistent state is made of. That is a separate question, and it deserves its own preparation before an answer is given.

\section{From Graphs to Fields}

Suppose one accepts that a persistent world-state is necessary. The next question is what kind of mathematical object it should be. The most familiar candidates all share a common commitment, and that commitment is worth making explicit before it is abandoned.

A scene graph represents a world as a discrete collection of entities connected by discrete relations: this object is on top of that object, this character is inside that room. A knowledge graph represents a world as a discrete collection of typed nodes connected by discrete typed edges: this fact holds of this entity. A voxel grid represents a world as a discrete lattice of cells, each holding some fixed local content. A polygon mesh represents a world as a discrete collection of vertices, edges, and faces approximating a continuous surface. Each of these representations differs enormously in what it is good for, yet each begins from the same assumption: that the world is fundamentally a collection of discrete, individuated things, and that representing the world means enumerating those things and their relations.

This assumption is not wrong so much as it is premature. Discreteness is a property that some regions of a well-behaved world exhibit, not a starting axiom that the representation should impose from the outset. A stone wall is, at a sufficiently coarse grain, well described as a discrete object with discrete properties. But the moss growing on it, the gradient of dampness near its base, the diffuse light falling across its surface, and the slow accumulation of soot from a nearby fire are not naturally discrete at all. Forcing them into a graph or a voxel grid means choosing a resolution and a set of categories in advance, and every choice of resolution and category is a choice about what the representation is permitted to notice. Discrete representations are efficient exactly to the degree that they have already decided what matters, and brittle exactly to the degree that something outside that decision turns out to matter after all.

A field-based representation inverts this commitment. Rather than beginning with entities and asking what properties they have, a field begins with a domain, typically a manifold of positions or a more abstract space of situations, and assigns to every point in that domain a value describing the local state of the world there. Objects, on this view, are not primitive. They are patterns that the field happens to exhibit: regions of high coherence, sharp gradients marking boundaries, stable configurations that persist under the field's own dynamics. Nothing is lost by allowing objects to emerge this way rather than positing them from the start, and a great deal is gained, because the same underlying representation can describe the wall, the moss, the dampness gradient, and the soot without deciding in advance which of these deserves to be a first-class entity and which is mere texture.

A careful reader might object that this is a problem of resolution rather than kind: why not simply use a much finer graph, with enough nodes and continuous-valued attributes to approximate whatever granularity the moss or the dampness requires? The objection misses what the discreteness commitment actually costs. However fine a graph becomes, building it still requires deciding, before a single value is written down, what counts as a node. Fineness does not remove that decision; it only postpones the point at which it has to be made explicit, and it multiplies the number of times it has to be made. A field requires no such decision. It assigns a value to every point of its domain from the outset, and whatever nodes, boundaries, or objects a graph would have needed to posit in advance become things one can instead read off the field's own structure: a boundary is simply where the field's values change sharply, an object is simply a region where they cohere. The field does not do without structure. It postpones commitment to structure until the structure has actually revealed itself, and that is precisely the property a persistent world-state needs.

This is the motivation for treating the persistent world-state as a field rather than a graph, and it is the motivation this book will build on for the remainder of the chapter.

\section{Defining the Semantic Field}

With that motivation in place, the central object of this book can now be stated. The world-state at time t is a field over a domain X, taking values in some appropriate space:

\[ M_t : X \to \mathbb{R}^d. \]

Here X is the domain over which the world is defined. Depending on the application, X might be a literal spatial manifold, a more abstract semantic space, or some hybrid of the two; the formalism does not require committing to one interpretation at this stage. What matters is that \(M_t\) assigns, to every point of X, a vector of real values describing the local state of the world at that point and at that moment. Where a scene graph would ask whether a wall-object exists and what properties it has, the field simply reports a value at every point of the domain, and wall-like coherence is something one can observe in how those values are organized, not something one has to declare up front.

This single equation is the culmination the chapter has been building toward, and it is worth pausing on the fact that so much preparation was required before it could be written down honestly. Introduced without the preceding sections, \(M_t\) : X → \(\mathbb{R}^d\) looks like an arbitrary formal choice, one field-theoretic representation among many that could have been picked for reasons of mathematical convenience. Introduced after the preceding sections, it looks like the answer to a question the reader has already been made to ask: if the wall must persist between observations, and if persistence should not be purchased by forcing the world into a premature discreteness, what is left? A field over a domain, reporting local state, from which observations are drawn by projection rather than invention. That is what \(M_t\) is for.

\section{Observation Operators}

It is worth pausing to name explicitly something that section 6.2 already introduced in passing, because it will recur throughout this book in increasingly sophisticated forms. The operator \(O_\theta\), mapping the field \(M_t\) to an observation \(y_t\), is not a detail of implementation. It is the second primitive this book depends on, standing beside the field itself.

The field and the observation operator answer two different questions, and keeping them separate matters more than it might first appear. The field answers what exists: \(M_t\) persists whether or not anything is currently reading it, in the same sense that the stone wall persists on a night when no one walks past it. The observation operator answers what is available to a particular observer at a particular moment, through a particular means of access: a human eye, a camera, a language description, a sensor array. Two observers with different operators, reading the same field, may recover different aspects of it, or recover the same aspects at different fidelity, without this implying that the field itself is different or that either observer is mistaken. Nothing about \(O_\theta\) constructs the wall. It only ever reports on a wall that was already there to be reported on.

This chapter will say nothing further about what \(O_\theta\) looks like concretely, and that restraint is deliberate. The question of which quantities a given observation can actually recover, and the question of how something like a camera or a renderer might realize \(O_\theta\) in practice, both depend on the field's internal structure, which has not yet been introduced. What matters here is only that the operator exists, that it is distinct from the field it reads, and that the distinction between world and observer will turn out to be one of this book's organizing threads rather than a footnote.

It is possible, even at this early stage, to sketch the complete loop that observation belongs to, without yet giving any of its parts a formal treatment. The field is observed, producing \(y_t\). Some process of reasoning, not yet specified, turns that observation into a proposed change, written \(\Delta\psi\). That proposal is not simply accepted; it is projected onto whatever states the field's own constraints permit, an operation this book will later write as \(\Pi_C\), and only the result of that projection is committed, becoming the field's next state \(M_{t+1}\):

\[ M_t \xrightarrow{O_\theta} y_t \xrightarrow{\text{reasoning}} \Delta\psi \xrightarrow{\Pi_C} Y \xrightarrow{\text{Commit}} M_{t+1}. \]

Every box in that diagram will be revisited. Part III gives \(\Delta\psi\), \(\Pi_C\), and Commit their full treatment as distinct operators rather than a single undifferentiated update. Part IV supplies the mathematics that makes projection well defined. Part V shows how each stage is actually implemented. Nothing about the diagram will change as the book proceeds; what changes is only how much is written inside each box. It is worth carrying this loop forward from here, because nearly everything the book develops from this point on is an elaboration of one of its five stages.

\section{Why Fields Instead of Objects}

Treating the world as a field rather than a collection of objects changes what identity means. In an object-based ontology, identity is typically a label: this entity is the same entity across time because it carries the same identifier, the same handle, the same node in the graph. Labels are cheap to assign and expensive to justify. Nothing about attaching the same label to two states of affairs guarantees that anything actually persisted between them; the label can simply be reasserted by fiat.

In a field-based ontology, identity has to be earned differently. An object is not a label attached to a region of the field. It is a stable configuration that the field's own dynamics maintain over time: a coherent, low-entropy, self-reinforcing pattern that persists because the field's evolution keeps reproducing it, not because something outside the field decided to keep calling it by the same name. The wall is a wall not because a symbol says so but because the relevant region of the field remains coherent, structurally continuous, and dynamically stable across the interval between two observations. If that region were to dissolve, drift, or reorganize beyond recognition, no label could rescue its identity; the wall would, in the most literal sense available to this framework, have stopped being the same wall.

This reframing has a further consequence that will not be developed fully until much later in the book, but that deserves to be planted here while the ground is fresh. If identity is the persistence of a trajectory through the field's state space rather than the persistence of a label, then questions about what a thing is become questions about what has continued, rather than questions about what has been asserted. This is a process ontology in the sense philosophers of that name would recognize, though this book has no need to argue for that tradition on its own terms. It is enough, for now, to notice that treating the world as a field rather than a collection of labeled objects was never merely a modeling convenience. It changes what it means for something to be the same thing twice, which is precisely the property that a stone wall needed and that a naive generative system could not supply.

\section{Looking Ahead}

This chapter has deliberately withheld an answer. \(M_t\) : X → \(\mathbb{R}^d\) says that the world is a field and that every point of that field carries a vector of real values, but it says nothing about what those values encode. A vector in \(\mathbb{R}^d\) is silent about whether it represents color, material, uncertainty, history, or something else again. That silence was appropriate here: the chapter's task was to establish that a persistent field is the right kind of object for a world-state to be, not to specify its contents.

The next chapter takes up that specification directly. If the world is a persistent field, what does each point in that field actually contain? That is the one question this chapter leaves open, and it is the question Chapter 7 exists to answer.
